\documentclass[../../../include/open-logic-section]{subfiles}

\begin{document}
\olfileid{sth}{card-arithmetic}{expotough}

\olsection[થોડાં સરળીકરણો]{ગણાંક ઘાતક્રિયાનાં થોડાં સરળીકરણો}

$\canonord$ની વ્યાખ્યા થોડી જટિલ હતી, પરંતુ તેનાથી
ગણાંક સરવાળા અને ગુણાકાર માટે ઉપયોગી
\olref[simp]{cardplustimesmax} મળ્યું. જોકે અનંતાતીત
ઘાતક્રિયાનું સરળીકરણ એટલી સહેલાઈથી થતું નથી.
શા માટે તે સમજાવવા સાન્ત કિસ્સાની પરિચિત રીતને
વિસ્તારતું પરિણામ લઈએ (તેનો પુરાવો ઊંચા અમૂર્ત
સ્તરે છે):

\begin{prop}\ollabel{simplecardexpo}
કોઈ પણ ગણાંકો $\cardfont{a}, \cardfont{b},
\cardfont{c}$ માટે
$\cardexpo{\cardfont{a}}{\cardfont{b} \cardplus \cardfont{c}} =
\cardexpo{\cardfont{a}}{\cardfont{b}} \cardtimes
\cardexpo{\cardfont{a}}{\cardfont{c}}$ અને
$\cardexpo{(\cardexpo{\cardfont{a}}{\cardfont{b}})}{\cardfont{c}} =
\cardexpo{\cardfont{a}}{\cardfont{b} \cardtimes \cardfont{c}}$.
\end{prop}

\begin{proof}
પહેલા દાવા માટે $f \colon
(\cardfont{b}\disjointsum\cardfont{c}) \to
\cardfont{a}$ વિધેય લો. તેને ``બે ભાગમાં વહેંચો'':
દરેક $\beta \in \cardfont{b}$ માટે
$f_\cardfont{b}(\beta) = f(\beta, 0)$ અને દરેક
$\gamma \in \cardfont{c}$ માટે
$f_\cardfont{c}(\gamma) = f(\gamma, 1)$
વ્યાખ્યાયિત કરો. નકશો $f \mapsto
\tuple{f_{\cardfont{b}}, f_\cardfont{c}}$ એ
!!a{bijection}
$\funfromto{\cardfont{b} \disjointsum \cardfont{c}}{\cardfont{a}} \to
(\funfromto{\cardfont{b}}{\cardfont{a}} \times
\funfromto{\cardfont{c}}{\cardfont{a}})$ છે.
\sourcecorrection{OLMD-137}{મૂળમાં અહીં નકશાનું મૂલ્ય
$(f_{\cardfont{b}} \times f_\cardfont{c})$
લખાયું છે. તે બે વિધેયોની કાર્ટેસિયન-ગુણાકાર
ક્રિયા સૂચવે છે, જ્યારે દર્શાવેલા લક્ષ્ય-ગણનો
ઘટક વિધેયોની ક્રમિત જોડી છે. તેથી અહીં
$\tuple{f_{\cardfont{b}}, f_\cardfont{c}}$ લખ્યું છે.}

બીજા દાવા માટે $f \colon \cardfont{c} \to
(\funfromto{\cardfont{b}}{\cardfont{a}})$ વિધેય લો;
એટલે દરેક $\gamma \in \cardfont{c}$ માટે
$f(\gamma) \colon \cardfont{b} \to \cardfont{a}$
એવું વિધેય છે. હવે દરેક
$\tuple{\beta, \gamma} \in \cardfont{b} \times
\cardfont{c}$ માટે
$f^*(\beta, \gamma) = (f(\gamma))(\beta)$
વ્યાખ્યાયિત કરો. $f \mapsto f^*$ એ
!!a{bijection}
$\funfromto{\cardfont{c}}{(\funfromto{\cardfont{b}}{\cardfont{a}})}
\to \funfromto{\cardfont{b} \times \cardfont{c}}{\cardfont{a}}$ છે.
\sourcecorrection{OLMD-138}{મૂળમાં આ નકશાનું
લક્ષ્ય-ગણ $\funfromto{\cardfont{b} \cardtimes
\cardfont{c}}{\cardfont{a}}$ લખાયું છે, પરંતુ
અહીં રચેલું $f^*$ સીધું
$\cardfont{b}\times\cardfont{c}$ પર વ્યાખ્યાયિત છે.
એ ગુણાકાર-ગણનો ગણાંક
$\cardfont{b}\cardtimes\cardfont{c}$ છે;
આ બે પ્રદેશ વચ્ચે એક-એક અને વ્યાપ્ત વિધેયથી
સ્થાનાંતરણ કરીને મૂળનું ગણાંક-સમીકરણ મળે છે.}
\end{proof}

અનંત ગણાંકો સાથે કામ કરીએ ત્યારે
$\cardexpo{\cardfont{a}}{\cardfont{b}}$ સહેલાઈથી
ગણી શકીએ એવી રીત આપણને \emph{ગમશે}.
આ દિશામાં એક સારું પગલું આ છે:

\begin{prop}\ollabel{cardexpo2reduct}
જો $2 \leq \cardfont{a} \leq \cardfont{b}$ અને
$\cardfont{b}$ અનંત હોય, તો
$\cardexpo{\cardfont{a}}{\cardfont{b}} =
\cardexpo{2}{\cardfont{b}}$.
\end{prop}

\begin{proof}
\begin{align*}
	\cardexpo{2}{\cardfont{b}} &\leq 
	\cardexpo{\cardfont{a}}{\cardfont{b}}\text{, કારણ કે $2 \leq \cardfont{a}$}\\
	&\leq \cardexpo{(2^\cardfont{a})}{\cardfont{b}}
	\text{, \olref[opps]{lem:SizePowerset2Exp} મુજબ}\\
	&= \cardexpo{2}{\cardfont{a} \cardtimes \cardfont{b}}
	\text{, \olref{simplecardexpo} મુજબ} \\
	&= \cardexpo{2}{\cardfont{b}}
	\text{, \olref[simp]{cardplustimesmax} મુજબ}
\end{align*}
\end{proof}

ખરેખર આનાથી વધુ સરળીકરણની અપેક્ષા રાખવી
જોઈએ નહીં, કારણ કે
\olref[card-arithmetic][opps]{lem:SizePowerset2Exp}
મુજબ $\cardfont{b} < \cardexpo{2}{\cardfont{b}}$.
પરંતુ $\cardfont{b} < \cardfont{a}$ હોય ત્યારે
$\cardexpo{\cardfont{a}}{\cardfont{b}}$ વિશે શું
કહેવું તે આથી ખબર પડતી નથી. અલબત્ત,
$\cardfont{b}$ \emph{સાન્ત} હોય તો શું કરવું તે જાણીએ છીએ.

\begin{prop}
$\cardfont{a}$ અનંત હોય અને $0<n \in \omega$ હોય, તો
$\cardexpo{\cardfont{a}}{n} = \cardfont{a}$.
\sourcecorrection{OLMD-139}{મૂળમાં $n \in \omega$
જ લખાયું છે; પરંતુ $n=0$ હોય ત્યારે
$\cardexpo{\cardfont{a}}{0}=1$, જે અનંત
$\cardfont{a}$ જેટલું નથી. તેથી અહીં $n>0$
શરત દૃશ્ય રીતે ઉમેરેલી છે.}
\end{prop}

\begin{proof}
$\cardexpo{\cardfont{a}}{n} = \cardfont{a} \cardtimes \cardfont{a}
\cardtimes \ldots \cardtimes \cardfont{a} = \cardfont{a}$,
\olref[simp]{cardplustimesmax} મુજબ.
\end{proof}
\noindent
વળી કેટલાક બીજા કિસ્સાઓમાં પણ
$\cardexpo{\cardfont{a}}{\cardfont{b}}$નું કદ
નિયંત્રિત કરી શકીએ છીએ:

\begin{prop}
જો $2 \leq \cardfont{b} < \cardfont{a} \leq
\cardexpo{2}{\cardfont{b}}$ અને $\cardfont{b}$ અનંત હોય, તો
$\cardexpo{\cardfont{a}}{\cardfont{b}} = \cardexpo{2}{\cardfont{b}}$.
\end{prop}

\begin{proof}
$\cardexpo{2}{\cardfont{b}}\leq \cardexpo{\cardfont{a}}{\cardfont{b}}
\leq \cardexpo{(\cardexpo{2}{\cardfont{b}})}{\cardfont{b}} =
\cardexpo{2}{\cardfont{b}\cardtimes\cardfont{b}} =
\cardexpo{2}{\cardfont{b}}$;
\olref{cardexpo2reduct} જેવી દલીલ છે.
\end{proof}
\noindent
પરંતુ આથી આગળ વાત ઘણી વધુ સૂક્ષ્મ બને છે.

\end{document}
